%%
%% main.tex — MatinBook Template
%% 
%% This is the main entry point for a MatinBook book.
%% It demonstrates the recommended structure and order of a
%% Persian technical book according to Iranian publishing standards.
%%
%% IMPORTANT:
%%   - This template uses \makecover (TikZ), which requires
%%     AT LEAST 2 compilation passes.
%%   - For a full build (with bibliography and index), use:
%%       xelatex -shell-escape main.tex
%%       biber main
%%       xelatex -shell-escape main.tex
%%       xelatex -shell-escape main.tex
%%     Or use latexmk:
%%       latexmk -xelatex -shell-escape main.tex
%%

\documentclass{matinbook}

%========================================================================
% BIBLIOGRAPHY
%========================================================================
\addbibresource{references.bib}

%========================================================================
% BOOK METADATA
%========================================================================
% These are used by \maketitle. \makecover uses its own arguments.
% Keep them in sync to avoid inconsistencies.
\title{عنوان کتاب}
\author{نام نویسنده}
\date{تابستان ۱۴۰۴}

% Repository (shown on cover)
\renewcommand{\repository}{github.com/your-username/your-book}

%========================================================================
% DOCUMENT
%========================================================================
\begin{document}

%------------------------------------------------------------------------
% FRONT COVER
%------------------------------------------------------------------------
% \makecover{title}{subtitle}{author}{date}
\makecover
    {عنوان کتاب}
    {زیرعنوان کتاب}
    {نام نویسنده}
    {تابستان ۱۴۰۴}

%------------------------------------------------------------------------
% FRONT MATTER (Persian letter numbering: الف، ب، ج، ...)
%------------------------------------------------------------------------
\frontmatter

% Title page (uses \title, \author, \date)
\maketitle

% Copyright page
\thispagestyle{empty}
\vfill
\begin{center}
    {\large\textbf{حقوق نشر}}
    
    \vspace{1cm}
    تمامی حقوق این کتاب محفوظ است.
    
    \vspace{1cm}
    \textbf{شناسنامه کتاب}
    
    \vspace{0.5cm}
    \begin{tabular}{rl}
        عنوان: & عنوان کتاب \\
        نویسنده: & نام نویسنده \\
        ناشر: & نام ناشر \\
        نوبت چاپ: & اول \\
        تیراژ: & ۱۰۰۰ نسخه \\
        شابک: & ۹۷۸-XXX-XXX-XXX-X \\
        قیمت: & XXX,XXX ریال \\
    \end{tabular}
    
    \vspace{1cm}
    {\small نسخه \matinbookversion\ — \matinbookdate}
\end{center}
\clearpage

% Dedication (optional)
\thispagestyle{empty}
\vfill
\begin{flushright}
    {\Large\itshape
        تقدیم به...
    }
\end{flushright}
\clearpage

% Bismillah page (optional)
% Uncomment to enable:
% \thispagestyle{empty}
% \vfill
% \begin{center}
%     {\Huge\bfseries بِسْمِ اللَّهِ الرَّحْمَٰنِ الرَّحِيمِ}
% \end{center}
% \vfill
% \clearpage

% Preface
\chapter*{پیشگفتار}
\addcontentsline{toc}{chapter}{پیشگفتار}

پیشگفتار کتاب را اینجا بنویسید...

\clearpage

% Table of contents
\tableofcontents
\clearpage

% List of figures
\listoffigures
\clearpage

% List of tables
\listoftables
\clearpage

% List of algorithms
% Uncomment ONLY if you use algorithm environments:
% \listofalgorithms
% \clearpage

%------------------------------------------------------------------------
% MAIN MATTER (Arabic numbering: 1, 2, 3, ...)
%------------------------------------------------------------------------
\mainmatter

% Chapters
% \input{chapters/chapter-01}
% \input{chapters/chapter-02}
% \input{chapters/chapter-03}
% \input{chapters/chapter-04}
% \input{chapters/chapter-05}
% \input{chapters/chapter-06}
% \input{chapters/chapter-07}
% \input{chapters/chapter-08}

%------------------------------------------------------------------------
% BACK MATTER
%------------------------------------------------------------------------
\backmatter

% Bibliography
\printbibliography[title={منابع و مراجع}]

% Index
\printindex

% Latin title page (optional)
% Uncomment to enable:
% \clearpage
% \thispagestyle{empty}
% \begin{latin}
% \begin{center}
%     \vspace*{3cm}
%     {\Huge\bfseries \latintitle} \\
%     \vspace{1cm}
%     {\Large \latinauthor}
% \end{center}
% \end{latin}

%------------------------------------------------------------------------
% BACK COVER (must be at the very end)
%------------------------------------------------------------------------
\makebackcover{%
    توضیحات پشت جلد کتاب را اینجا بنویسید...
}

\end{document}
